% ===========================================================================
%  E01 — Entrar e circular na plataforma
% ===========================================================================
\chapter{E01 \textperiodcentered{} Entrar e circular na plataforma}
\label{ep:e01}

\begin{fichaepico}
\begin{campos}
  \item[Objetivo] Cada pessoa entra com uma conta só, cai direto na área do
    próprio papel e circula sem se perder --- inclusive quem trabalha em mais
    de uma escola.
  \item[Para quem] Todos os seis papéis com conta.
  \item[Valor] Às 8h da manhã, milhares de crianças entram ao mesmo tempo. Cada
    segundo de atrito na entrada é tempo de aula perdido.
  \item[Estado] \parcial{} --- entrar, primeiro acesso e navegação estão
    prontos; faltam a recuperação de senha e o login único da rede, e a
    Ajuda com atendente está construída, em teste interno.
\end{campos}
\end{fichaepico}

\begin{tabularx}{\linewidth}{@{}L{16mm}X L{20mm}@{}}
\toprule
\headrow \thd{ID} & \thd{Feature} & \thd{Estado}\\
\midrule
\id{F01.1} & Entrar com usuário e senha & \pronto\\
\zebra \id{F01.2} & Cada pessoa na sua área & \pronto\\
\id{F01.3} & Trabalhar em mais de uma escola & \pronto\\
\zebra \id{F01.4} & Cuidar da própria senha & \parcial\\
\id{F01.5} & Entrar pelo login único da rede & \afazer\\
\zebra \id{F01.6} & Tirar dúvidas com o atendente da Ajuda & \parcial\\
\bottomrule
\end{tabularx}

% ---------------------------------------------------------------------------
\feature{F01.1}{Entrar com usuário e senha}{\pronto}

Uma única tela de entrada serve a todos os papéis. É o papel da pessoa que
decide para onde ela vai depois.

\begin{historia}{H01.1.1}{Entrar e cair na minha área}{\pronto}
\hu{qualquer pessoa com conta na rede}{entrar com usuário e senha e cair
direto na minha área}{não ter de procurar por onde começar}
\aceite
\begin{itemize}
  \item A tela se chama \enquote{Entrar na plataforma}, com os campos Usuário
    e Senha e o botão \enquote{Entrar}.
  \item Depois de entrar: o Estudante vai para \enquote{Minhas missões}; o
    Professor para o Início dele; as três coordenações e o Administrador do
    contrato para \enquote{Usuários}; o Administrador Unooria para
    \enquote{Recursos}.
  \item Senha errada, usuário inexistente e conta desativada recebem
    \emph{a mesma} mensagem: \enquote{Usuário ou senha incorretos.} --- a tela
    nunca diz qual das duas metades errou, para não confirmar que um usuário
    existe.
  \item Depois de muitas tentativas seguidas, aparece \enquote{Muitas
    tentativas. Aguarde um momento e tente de novo.}
  \item Se a plataforma estiver com problema, aparece \enquote{Não foi possível
    entrar agora. Tente de novo em alguns instantes.}
\end{itemize}
\end{historia}

\begin{historia}{H01.1.2}{Criar minha senha no primeiro acesso}{\pronto}
\hu{pessoa que recebeu uma senha provisória (da importação ou de um
reajuste)}{ser levada a criar a minha própria senha antes de qualquer
coisa}{que só eu conheça a minha senha}
\aceite
\begin{itemize}
  \item A tela \enquote{Crie sua senha} aparece logo depois de entrar, sem menu
    e sem saída.
  \item Campos: \enquote{Senha atual}, \enquote{Nova senha} (pelo menos 8
    caracteres, diferente da atual) e \enquote{Repita a nova senha}.
  \item Enquanto a senha não for trocada, nada funciona: nem as telas nem as
    missões. Digitar o endereço de uma missão não pula esta etapa.
  \item Vale para todos os papéis, inclusive o Administrador Unooria.
  \item Trocada a senha, a pessoa segue para a área dela normalmente.
\end{itemize}
\end{historia}

\begin{historia}{H01.1.3}{Não cair no meio da aula}{\pronto}
\hu{Estudante}{que a minha sessão não expire enquanto estou numa
missão}{não perder a aula tendo de entrar de novo}
\aceite
\begin{itemize}
  \item A sessão expira depois de 2 horas \textbf{sem uso}, igual para todos os
    papéis.
  \item Ficar dentro de uma missão conta como uso.
\end{itemize}
\end{historia}

\begin{historia}{H01.1.4}{Sair}{\pronto}
\hu{qualquer pessoa}{sair da plataforma}{que ninguém use a minha conta num
computador compartilhado da escola}
\aceite
\begin{itemize}
  \item O menu do avatar, no cabeçalho, tem o item \enquote{Sair}.
  \item Depois de sair, voltar a qualquer tela interna leva ao login.
\end{itemize}
\end{historia}

% ---------------------------------------------------------------------------
\feature{F01.2}{Cada pessoa na sua área}{\pronto}

\begin{historia}{H01.2.1}{Voltar para a minha área sem ver erro}{\pronto}
\hu{Estudante que digitou o endereço da área do professor}{ser levado de
volta à minha área, sem mensagem de erro}{não ficar perdido nem assustado}
\aceite
\begin{itemize}
  \item Quem entra no endereço de outra área volta para a própria, sem aviso de
    \enquote{sem permissão}.
  \item A proteção não depende disso: os dados de outra área também são
    recusados, mesmo que alguém tente por outro caminho.
\end{itemize}
\end{historia}

\begin{historia}{H01.2.2}{Saber onde estou}{\pronto}
\hu{qualquer pessoa}{ver onde estou e voltar um nível com um clique}{não me
perder numa plataforma com muitas telas}
\aceite
\begin{itemize}
  \item O cabeçalho é o mesmo em todas as áreas; a marca leva ao início do
    papel.
  \item Toda tela interna tem a trilha de navegação (\enquote{Início \seta{}
    Entre Redes \seta{} Missão\ldots}), exceto o início de cada papel.
  \item Dentro de uma missão, o cabeçalho mostra \enquote{\ensuremath{\leftarrow}
    Missões}.
  \item O cabeçalho mostra nome, uma linha de apoio (a série do estudante ou o
    papel) e um avatar com as iniciais.
  \item As telas de gestão têm menu lateral filtrado pelo papel.
  \item Os menus laterais da gestão e da Unooria mostram só o que leva a algum
    lugar: o item \enquote{Iniciativas}, que não fazia nada, saiu.
\end{itemize}
\end{historia}

% ---------------------------------------------------------------------------
\feature{F01.3}{Trabalhar em mais de uma escola}{\pronto}

\begin{historia}{H01.3.1}{Trocar a escola ativa}{\pronto}
\hu{Professor ou Coordenação com vínculo em mais de uma escola}{trocar a
escola em que estou trabalhando pelo cabeçalho}{ver e registrar só o que é
daquela escola}
\aceite
\begin{itemize}
  \item Com vínculo em 2 ou mais escolas aparece o seletor \enquote{Escola
    ativa}; com uma escola só, ele não aparece.
  \item Trocar de escola muda tudo o que a pessoa vê, sem recarregar a página.
  \item As ferramentas abertas (plano de aula, planejamento) recomeçam na
    escola nova, para nada ser gravado na escola errada.
  \item No celular o seletor continua no cabeçalho, mais estreito; tocá-lo
    abre a lista do próprio telefone, com o nome inteiro de cada escola.
\end{itemize}
\end{historia}

% ---------------------------------------------------------------------------
\feature{F01.4}{Cuidar da própria senha}{\parcial}

Hoje a senha só muda quando uma coordenação usa \enquote{Resetar senha}
(\id{H03.3.4}). O item \enquote{Trocar senha} ficou fora do menu de propósito,
para não ser um link que não leva a nada.

\begin{historia}{H01.4.1}{Trocar minha senha quando eu quiser}{\pronto}
\hu{qualquer pessoa}{trocar a minha senha a qualquer momento}{me proteger se
alguém descobriu a senha}
\aceite
\begin{itemize}
  \item O menu do avatar tem \enquote{Trocar senha}, com as mesmas regras do
    primeiro acesso: a senha atual, pelo menos 8 caracteres e diferente da
    atual.
  \item A tela se chama \enquote{Trocar senha} e tem \enquote{Cancelar}, que
    volta para a área da pessoa --- ela não fica presa ali, como no primeiro
    acesso.
  \item Depois de salvar, a tela confirma \enquote{Senha trocada} e
    \enquote{Voltar} leva de volta à área da pessoa. Na próxima entrada vale a
    senha nova.
\end{itemize}
\end{historia}

\begin{historia}{H01.4.2}{Recuperar a senha esquecida}{\afazer}
\hu{Professor ou Coordenação}{recuperar a minha senha sem depender de
alguém}{não ficar sem acesso no dia da aula}
\aceiteproposta
\begin{itemize}
  \item Na tela de entrada, \enquote{Esqueci minha senha} envia um link de uso
    único para o e-mail cadastrado.
  \item Para o estudante (que muitas vezes não tem e-mail), o caminho continua
    sendo a coordenação --- ver a decisão \id{D05} no
    Capítulo~\ref{cap:decisoes}.
\end{itemize}
\end{historia}

% ---------------------------------------------------------------------------
\feature{F01.5}{Entrar pelo login único da rede}{\afazer}

A rede usa um provedor de identidade central (WSO2), o mesmo dos outros
sistemas da secretaria. A plataforma já está preparada para ele: a troca é de
configuração. Falta validar num ambiente de homologação com a rede.

\begin{historia}{H01.5.1}{Usar a conta que eu já tenho na rede}{\afazer}
\hu{qualquer pessoa da rede}{entrar com a mesma conta que uso nos outros
sistemas da secretaria}{não ter mais uma senha para decorar}
\aceiteproposta
\begin{itemize}
  \item A tela de entrada leva ao login da rede; ao voltar, a pessoa cai na
    mesma área de hoje.
  \item Nada muda dentro das missões nem nas telas.
  \item O primeiro acesso com senha provisória deixa de existir para quem
    entra pelo login da rede.
\end{itemize}
\end{historia}

% ---------------------------------------------------------------------------
\feature{F01.6}{Tirar dúvidas com o atendente da Ajuda}{\parcial}

Um atendente que conhece a plataforma. A pessoa abre
\enquote{Ajuda} no menu do próprio nome, no alto da tela, e pergunta falando
ou escrevendo; o atendente responde em voz e o texto aparece embaixo dele.
São dois atendentes: o \textbf{mascote da Unooria}, um personagem 3D que fala
com a pessoa (\id{H01.6.4}), e o \textbf{atendente em vídeo}, uma pessoa
virtual de um serviço contratado. Quem atende primeiro é escolha do
Administrador Unooria (\id{H10.4.3}); hoje, o mascote.
Ele sabe o que este mapa diz que está pronto, os guias da aula e as missões
--- e mais nada: não vê turma, estudante nem resultado.

\begin{nota}
\textbf{O tempo de conversa é limitado.} Cada conversa e cada dia têm um teto
de minutos, que o Administrador Unooria define (\id{H10.4.2}). Começa em 10
minutos por conversa e 30 por dia, por pessoa. A Ajuda não gasta o limite
diário de IA de ninguém (\id{F09.4}). Os minutos são os mesmos para os dois
atendentes: trocar de um para o outro não dá tempo novo.
\end{nota}

\begin{historia}{H01.6.1}{Perguntar ao atendente sem sair da tela}{\parcial}
\hu{Professor, Coordenação ou Administrador}{abrir a Ajuda de qualquer tela e
perguntar ao atendente como fazer alguma coisa na plataforma}{resolver a
dúvida na hora, sem esperar alguém do suporte}
\aceite
\begin{itemize}
  \item \enquote{Ajuda} é o segundo item do menu do nome, para todos os papéis
    menos o Estudante. Abre uma janela por cima da tela: o vídeo do atendente
    em cima, a conversa escrita embaixo.
  \item Enquanto conecta, a janela diz \enquote{Chamando o atendente\ldots} e
    depois \enquote{Conectando o vídeo\ldots}; o atendente se apresenta e
    espera a pergunta.
  \item O atendente responde sobre o que a plataforma faz hoje, no papel de
    quem pergunta. O que não existe ainda, ele diz que não existe.
  \item Junto do campo, fixo: \enquote{Não diga nem escreva nomes de
    estudantes.} Nada da conversa é guardado: fechar a janela encerra a
    conversa e o vídeo.
  \item Se o vídeo não conectar, a janela diz que a rede da escola pode estar
    bloqueando a conexão e oferece \enquote{Tentar de novo}.
\end{itemize}
\pendencia{a Ajuda nasce desligada em toda rede, enquanto está em teste
interno: só aparece onde o Administrador Unooria a ligou (\id{H10.1.5}).
Falta o teste de ponta a ponta, e a voz do atendente fala português com
sotaque --- aceito pelo dono do produto.}
\end{historia}

\begin{historia}{H01.6.2}{Escrever quando não posso falar}{\parcial}
\hu{Professor ou Coordenação}{escrever a minha pergunta em vez de falar}{usar
a Ajuda numa sala com gente, ou num computador sem microfone}
\aceite
\begin{itemize}
  \item O microfone pode ser silenciado e religado a qualquer momento; o
    estado aparece no botão (\enquote{Silenciar microfone},
    \enquote{Ligar microfone}, \enquote{Microfone bloqueado}).
  \item Sem microfone, a janela diz \enquote{Sem microfone: escreva a sua
    pergunta abaixo.}, e a conversa segue só por escrito; o atendente responde
    em voz e por escrito do mesmo jeito.
  \item A pergunta escrita e a falada entram na mesma conversa, na ordem em
    que aconteceram.
\end{itemize}
\pendencia{a mesma de \id{H01.6.1}.}
\end{historia}

\begin{historia}{H01.6.3}{Saber quando o meu tempo de Ajuda acabou}{\parcial}
\hu{qualquer pessoa que usa a Ajuda}{saber quando a conversa acabou por
tempo, e quando o tempo do dia acabou}{não achar que o atendente travou}
\aceite
\begin{itemize}
  \item Ao chegar no limite da conversa, o atendente se despede, e a janela
    mostra \enquote{A conversa terminou}, com \enquote{Começar de novo}.
  \item Com os minutos do dia gastos, a Ajuda não abre e diz \enquote{Você já
    usou todos os minutos de Ajuda de hoje.} No dia seguinte, o tempo volta.
  \item Quando a resposta demora, o atendente avisa (\enquote{Um momento,
    estou vendo isso para você.}) em vez de ficar em silêncio; se não
    conseguir responder, pede que a pessoa repita a pergunta.
\end{itemize}
\pendencia{a mesma de \id{H01.6.1}.}
\end{historia}

\begin{historia}{H01.6.4}{Conversar com o mascote da Unooria}{\pronto}
\hu{Professor, Coordenação ou Administrador}{tirar a minha dúvida conversando
com o mascote da Unooria}{ser atendido por um personagem da própria
plataforma, que responde na hora, em voz e por escrito}
\aceite
\begin{itemize}
  \item Quando o mascote é quem atende, a janela da Ajuda mostra o mascote da
    Unooria em cima e a conversa escrita embaixo, no lugar do vídeo. Enquanto
    ele chega, a janela diz \enquote{Chamando o mascote\ldots} e depois
    \enquote{Conectando a conversa\ldots}.
  \item O mascote fala em voz alta, e a boca dele acompanha a fala. O que ele
    diz aparece por escrito na conversa, como no atendente em vídeo; o que a
    pessoa fala também.
  \item A pessoa pergunta falando ou escrevendo, do mesmo jeito que no
    atendente em vídeo: o microfone começa desligado e pode ser ligado e
    silenciado a qualquer momento; sem microfone, a conversa segue por escrito.
    Se a pessoa fala por cima, o mascote para e escuta.
  \item O mascote sabe o mesmo que o atendente em vídeo, responde no papel de
    quem pergunta e diz quando algo ainda não existe. Nada da conversa é
    guardado, e o aviso \enquote{Não diga nem escreva nomes de estudantes.}
    continua fixo junto do campo.
  \item Os minutos são os mesmos da Ajuda (\id{H01.6.3}): no fim do tempo da
    conversa, a janela mostra \enquote{A conversa terminou}, com
    \enquote{Começar de novo}.
  \item Se o Administrador Unooria deixou trocar de atendente, aparece um botão
    na janela --- \enquote{Conversar com o atendente em vídeo} ou
    \enquote{Conversar com o mascote} ---, e a conversa recomeça com o outro.
    Se o atendente escolhido não está disponível, a Ajuda abre com o outro;
    sem nenhum dos dois, diz \enquote{A Ajuda com atendente não está
    disponível agora.} e oferece \enquote{Tentar de novo}.
  \item Como toda a Ajuda, só aparece onde o Administrador Unooria a ligou
    (\id{H10.1.5}).
\end{itemize}
\end{historia}
